% paper/sections/11e_ao_fast_state_space.tex
%
% §11  AO-Fast geometric cell-fraction state space
%
% ═══════════════════════════════════════════════════════════════════════════════
\subsection{AO-Fast 幾何セル分率状態空間}
\label{sec:ao_fast_state_space}
% ═══════════════════════════════════════════════════════════════════════════════

本節は，ショートペーパー AO で定義した幾何セル分率状態空間を，
第2章の連続理論と第12--14章の実験 gate の間に置く数値アルゴリズム章である．
第~\ref{sec:capillary_pressure_reaction_subspace}節では，
表面エネルギー仮想仕事と圧力反力部分空間 $R_p$ を理論境界として固定した．
ここでは，その境界を高速に評価するための状態所有者，互換射影，
アクティブ集合，近似誤差，GPU 契約，および YAML/fail-close 契約を定式化する．

\subsubsection{状態所有者と幾何互換条件}
\label{sec:ao_state_ownership}

AO 経路では，物質量の所有者を平滑化された $\psi$ ではなく，
セルごとの有次元液相体積
\begin{equation}
  q_C = |C|\,\theta_C,\qquad
  \rho_C(q)=\rho_g+(\rho_l-\rho_g)\frac{q_C}{|C|}
  \label{eq:ao_volume_owner}
\end{equation}
に置く．
正規化分率 $\theta_C$ は表示・物性補間用の view であり，
非一様格子で $\sum_C\theta_C$ を体積として読んではならない．
鋭い毛管幾何は連続 gauge $\phi$ が誘導する界面複体
$\Gamma_h=\Gamma(\phi)$ が所有する．
固定正則 stratum $S$ 上では，液相領域
$\Omega_l(\phi)=\{\bx:\phi_h(\bx)<0\}$ に対して
\begin{equation}
  Q_h^S(\phi)_C
  =
  |C\cap\Omega_l(\Gamma(\phi))|,
  \qquad
  A_h^S(\phi)_C=\frac{Q_h^S(\phi)_C}{|C|},
  \qquad
  S_h^S(\phi)=|\Gamma(\phi)|
  \label{eq:ao_geometry_maps}
\end{equation}
を定義し，互換状態は
\begin{equation}
  Q_h^S(\phi)_C=q_C
  \quad (C\in\mathcal{A}_q)
  \label{eq:ao_geometric_compatibility}
\end{equation}
を満たす．
ここで $\mathcal{A}_q$ は混合セル，直前の混合セル，
輸送で状態が変わったセル，目標状態で $0<q_C<|C|$ となるセル，
および符号・境界所有権判定に必要な一面 halo からなるアクティブ制約支持である．
周期境界は商複体として数え，壁では phase flux が消える面を境界制約として扱う．
壁接触角を後から導入する場合は $dS_h$ の境界項で入れるべきであり，
曲率の後処理補正として混ぜてはならない．

この所有者分離により，拒否される経路も明確になる：
$q$ から密度を作りながら毛管力だけを互換でない $\psi$ から作る経路，
非幾何輸送後の大域質量補正，$q$ の clipping を保存性証明と呼ぶ経路，
および不連続な $\theta_C$ を CCD/DCCD/FCCD/UCCD の滑らかなスカラーとして微分する経路は，
AO 経路では採用しない．
CCD 族は速度，圧力，gauge predictor，smooth diagnostic，面状態再構成に使う．
幾何写像 $Q_h,J_q,T_q,dS_h$ の代替微分器としては使わない．

\subsubsection{互換射影と射影仕事}
\label{sec:ao_compatibility_projection}

保存輸送は，セル面 incidence $B$ と swept liquid flux $\Phi_l$ により
\begin{equation}
  q^{n+1}=q^n-\Delta t\,B\Phi_l,
  \qquad
  \mathbf{1}^T B\Phi_l=0
  \label{eq:ao_conservative_q_update}
\end{equation}
と書く．
密度・運動量の共通フラックスは同じ $\Phi_l$ から作り，
位相，密度，運動量を同一台帳で進める．
この段階で得た $q^-$ と gauge predictor $\phi^-$ は，互換 bundle
式~\eqref{eq:ao_geometric_compatibility} へ射影する：
\begin{equation}
  \min_{\phi}\frac12\|\phi-\phi^-\|_{W_\eta}^2
  \quad\text{subject to}\quad
  Q_h^S(\phi)_C=q^-_C\quad(C\in\mathcal{A}_q),
  \qquad
  W_\eta=W_v+\eta L^T H_L L .
  \label{eq:ao_projection_problem}
\end{equation}
線形化では $J_q=dQ_h^S/d\phi$ を用い，Schur 系
\begin{equation}
  S_q\lambda=J_q\delta_{\mathrm{pred}}-r_q,
  \qquad
  S_q=J_qW_\eta^{-1}J_q^T,
  \qquad
  \delta=\delta_{\mathrm{pred}}-W_\eta^{-1}J_q^T\lambda
  \label{eq:ao_projection_schur}
\end{equation}
を解く．
$J_q$ が宣言された圧力 gauge を除いて full row rank であれば，
$S_q$ は SPD であり，rank・条件数・sign/case margin・
物理体積残差 $\|Q_h(\phi^+)-q^-\|$ が射影受理条件になる．
この射影は gauge を互換状態へ戻す処理であり，
表面長の変化
\begin{equation}
  \Delta S_\Pi=S_h(\phi^+)-S_h(\phi^-)
  \label{eq:ao_projection_surface_work}
\end{equation}
は物理的な毛管仕事ではなく射影仕事として台帳に記録する．
大きすぎる $\Delta S_\Pi$ は fail-close gate であり，
表面張力の散逸や安定化として消費してはならない．

\subsubsection{毛管 bundle work と圧力反力分解}
\label{sec:ao_bundle_capillarity}

互換状態では，表面エネルギーは
$E_\sigma(\phi)=\sigma S_h(\phi)$ であり，
面速度 $w_f$ による体積変化は $T_q(\Gamma_h)w_f$ として読む．
この体積変化を互換 bundle に lift し，
その lift に対する $dS_h$ の作用を評価して，
面空間 $F_h$ 上の毛管共変量 $r_\sigma$ を得る．
したがって AO 経路の毛管力は曲率配列の後処理ではなく，
同じ $Q_h,J_q,T_q,dS_h,M_f$ が定める face-Hodge Riesz 代表である．

静的平衡は，あるセル圧力座標 $\pi$ が存在して
\begin{equation}
  \sigma\,dS_h[\delta\phi]+\pi^T J_q\delta\phi=0
  \quad
  \text{for all admissible }\delta\phi
  \label{eq:ao_discrete_young_laplace_test}
\end{equation}
を満たすかで判定する．
これは形状名による分岐ではなく，互換 bundle 上の Young--Laplace 条件である．
ただし，第~\ref{sec:capillary_pressure_reaction_subspace}節の
式~\eqref{eq:ao_full_pressure_cancellation} が示すように，
full cell-pressure image へ射影すると毛管波の非圧力駆動まで消える場合がある．
したがって実行時に速度を変える量は，
理論境界として先に定義した $R_p(q_T)$ に対する
式~\eqref{eq:capillary_pressure_reaction_projection} の
$M_f$-直交残差でなければならない．

\subsubsection{AO-Fast の計算量削減と近似誤差}
\label{sec:ao_fast_acceleration}

ショートペーパー AO の式を全領域で愚直に読むと，
各 time step と line search trial で全セルの cut geometry，
全セルの $J_q$ table，および全セル Schur 作用を再構築するため，
計算量は
\begin{equation}
  \text{direct AO:}\quad \Ord{k|C_h|}
  \quad\text{per candidate}
  \label{eq:ao_direct_complexity}
\end{equation}
となる．
AO-Fast の高速化単位は solver family の差し替えではなく，
実際に界面 stratum が変わったアクティブ表だけを更新することである：
\begin{equation}
  \text{AO-Fast:}\quad
  \Ord{|\mathrm{dirty}|+k|\mathcal{A}_q|}
  \label{eq:ao_fast_complexity}
\end{equation}
を目標とする．
ここで dirty は sign/case，交差区間，境界・周期所有権，目標支持，
swept-flux 支持，または metric identity が変化した行である．
通常実行時に $q^-$ を全領域 scan して mixed cell を探すことはしない．
transport が出す state-changing support stream と前 step の active table から
device 上で compaction し，full-grid support scan は初期化，restart 検証，
dense oracle 比較，明示 debug，または degenerate exact step に限定する．

固定正則 stratum では，凍結した局所幾何を proposal generator として使える．
局所 crossing margin と節点 sign margin から
\begin{equation}
  \beta_C
  =
  \frac{\|\delta\phi\|_{\infty,C}}
       {\min(\gamma_C,m_C)}
  \label{eq:ao_beta_margin}
\end{equation}
を定義し，$\beta_C\le\beta_\ast<1$ なら
\begin{align}
  |Q_h^S(\phi+\delta\phi)_C-Q_h^S(\phi)_C-J_{q,C}\delta\phi_C|
  &\le C_Q |C|\beta_C^2, \notag\\
  |S_h^S(\phi+\delta\phi)_C-S_h^S(\phi)_C-dS_C\delta\phi_C|
  &\le C_S |\Gamma_C|\beta_C^2 .
  \label{eq:ao_fixed_stratum_accuracy}
\end{align}
二次 secant/Hessian proposal を使う場合でも，
その精度主張は同じ stratum 上の候補精度であり，
commit は常に exact active recomputation の残差で判定する．

近似毛管分解の精度は，毛管加速度の仕事ノルムで宣言する．
近似共変量 $\widetilde r_\sigma$ と近似射影
$\widetilde\Pi_{R_p}^{M_f}$ が
\begin{equation}
  \|r_\sigma-\widetilde r_\sigma\|_{M_f^{-1}}\le\varepsilon_\sigma,\qquad
  \|\Pi_{R_p}^{M_f}r_\sigma-\widetilde\Pi_{R_p}^{M_f}\widetilde r_\sigma\|_{M_f^{-1}}
  \le\varepsilon_p
  \label{eq:ao_fast_error_budget}
\end{equation}
を満たすなら，
\begin{equation}
  \|a_{\sigma,\mathrm{bal}}-\widetilde a_{\sigma,\mathrm{bal}}\|_{M_f}
  \le \varepsilon_\sigma+\varepsilon_p
  \label{eq:ao_fast_acceleration_error_bound}
\end{equation}
である．
この境界は「近似を使ってよい」という許可ではなく，
近似を受理するために事前宣言された精度契約である．

\subsubsection{PCG/DC，GPU，YAML 契約}
\label{sec:ao_fast_solver_yaml_contract}

互換射影の主反復は，active graph 上の Schur 作用を matrix-free に評価する
PCG/Newton 系とし，前 step warm start，対角および connected-component block
preconditioner，必要なら component deflation を用いる．
反復許容誤差は downstream の exact gate で支配し，
\begin{equation}
  \tau_{\mathrm{cg}}
  \le
  \min(c_q\tau_q,\;c_s\tau_S,\;c_\kappa\tau_q/\kappa_{\mathrm{est}})
  \label{eq:ao_pcg_tolerance}
\end{equation}
のように，物理体積残差・射影仕事・条件数推定から決める．
丸め floor が目標を上回る場合は，さらに反復しても精度を買えないので，
fail-close または明示された solver-chain 遷移だけを許す．

DC は primary solver ではなく，exact residual を単調に下げる proposal としてのみ使う．
厳密 active residual を
\begin{equation}
  R(\phi)=Q_h^S(\phi)-q^-
  \label{eq:ao_exact_active_residual}
\end{equation}
とし，安価な凍結逆作用 $P_0$ から
\begin{equation}
  \delta\phi_{\mathrm{DC}}=-P_0R(\phi_k),
  \qquad
  \phi_{\mathrm{trial}}(\alpha)=\phi_k+\alpha\delta\phi_{\mathrm{DC}}
  \label{eq:ao_dc_candidate}
\end{equation}
を作る．
受理できるのは，
\begin{equation}
  \|R(\phi_{\mathrm{trial}})\|_{H_C^{-1}}
  <
  \|R(\phi_k)\|_{H_C^{-1}}
  \label{eq:ao_dc_residual_monotone}
\end{equation}
を exact recomputation で満たし，sign/case margin と
$\Delta S_\Pi$ gate も満たす場合だけである．
DC が停滞した場合に PCG へ移るかどうかは YAML で明示する方針であり，
暗黙 fallback ではない．

GPU 実装では，active table，cut-geometry row，$J/J^T$ 作用，
Krylov vector，line-search residual，rank/conditioning probe を
すべて \texttt{backend.xp} 上に保持する．
CG/Newton/DC/line-search の内側で \texttt{.get()}，\texttt{asnumpy}，
\texttt{float(...)}，\texttt{bool(...)}，Python list materialization による
host 同期を行わない．
host へ渡してよいのは，outer iteration 受理後の台帳スカラー，
kernel counter，host-transfer counter だけである．
dense direct-AO 由来の実装は，exact P1 case algebra，
metric cache，q/theta/phi 分離，fail-close parser gate，
face-Hodge work identity などの知見を取り込んでよいが，
本番経路では \texttt{oracle\_only}，\texttt{gpu\_production}，\texttt{reject}
の分類を通過した記号だけを呼ぶ．
全格子 dense projection は debug oracle であり，production fallback ではない．

YAML/UX では，primary solver，accelerator，fallback を分けて宣言する．
既定は fail-close であり，例えば
\texttt{fallback.policy: none} なら primary solver 失敗後に
PCG/DC を自動で切り替えない．
\texttt{fallback.policy: explicit\_chain} を選ぶ場合は，
\texttt{from}，\texttt{to}，\texttt{triggers}，\texttt{record\_as} を持つ
可視な chain として台帳に記録する．
同様に \texttt{geometric\_cell\_fraction} では，
輸送変数を \texttt{q}，互換制約を \texttt{hard\_cell\_volume} とし，
GPU 内側反復の host transfer と dense runtime fallback を禁止する．
\texttt{psi} 輸送，Ridge--Eikonal 再初期化，boundedness clipping，
diagnostic-only condition gate，implicit fallback は parse 時点で拒否する．

\subsubsection{実験章への接続}
\label{sec:ao_fast_validation_connection}

AO-Fast は，式~\eqref{eq:ao_geometric_compatibility} の hard volume constraint，
式~\eqref{eq:capillary_pressure_reaction_projection} の $R_p$ 分解，
式~\eqref{eq:ao_fast_error_budget}--\eqref{eq:ao_fast_acceleration_error_bound} の
近似誤差，および本節の GPU/YAML 契約を同時に満たすまでは，
第~\ref{sec:benchmarks}章の production 物理経路へ入れない．
第~\ref{sec:U12_ao_capillary_split_gate}節の U12 は full pressure image の反例と
GPU fail-close を単体 gate として確認し，
第~\ref{sec:ao_fast_capillary_gate}節の V11 は統合前 gate として
標準 FCCD/UCCD6/pressure-jump/component-Hodge 経路との境界を固定する．
したがって，第~\ref{sec:val_capillary}節の毛管波が完走しても，
それは AO-Fast production admission ではない．
AO-Fast の本番化は，active-vs-dense equality，dirty/flux/target refresh，
device-resident Schur/PCG，declared fallback ledger，
および Chapter 14 YAML の \texttt{geometric\_cell\_fraction} 明示化を
順に通過した後の別段階である．
